\documentclass[12pt,a4paper]{article}
\usepackage[utf8]{inputenc}
\usepackage[brazilian]{babel}
\usepackage[margin=2.5cm]{geometry}
\usepackage{amsmath}
\usepackage{listings}

\lstset{
    language=Python,
    basicstyle=\ttfamily\small,
    frame=single,
    numbers=left
}

\title{Transformações Geométricas 2D e 3D \\ \small EI 2 - Computação Gráfica}
\author{Nome do acadêmico}
\date{\today}

\begin{document}
\maketitle

\section*{Resumo}
Implementação em Python/Pygame de translação, escala e rotação em 2D e 3D. Dois programas: \texttt{transformacoes\_2d.py} (triângulo) e \texttt{transformacoes\_3d.py} (cubo com projeção isométrica).

\section{Fundamentação}
Um ponto $(x,y)$ em coordenadas homogêneas é $[x, y, 1]$. O ponto transformado é $P' = M \cdot P$.

\textbf{Translação:} $T(t_x, t_y)$ desloca o objeto.
\[
T = \begin{bmatrix} 1 & 0 & t_x \\ 0 & 1 & t_y \\ 0 & 0 & 1 \end{bmatrix}
\]

\textbf{Escala:} $S(s_x, s_y)$ altera o tamanho (em relação à origem).
\[
S = \begin{bmatrix} s_x & 0 & 0 \\ 0 & s_y & 0 \\ 0 & 0 & 1 \end{bmatrix}
\]

\textbf{Rotação:} $R(\theta)$ gira o objeto (ângulo em radianos).
\[
R = \begin{bmatrix} \cos\theta & -\sin\theta & 0 \\ \sin\theta & \cos\theta & 0 \\ 0 & 0 & 1 \end{bmatrix}
\]

Para escala ou rotação em torno de $(c_x, c_y)$: translação para o centro, aplicar S ou R, translação de volta.

\section{Implementação}
O programa usa um triângulo fixo no centro. Teclas: S (escala 1,2), R (rotação 15°), T (translação 30,20), Backspace (voltar).

\textbf{Translação} — soma $t_x$ e $t_y$ em cada coordenada:
\begin{lstlisting}
def translacao(pontos, tx, ty):
    return [(x + tx, y + ty) for x, y in pontos]
\end{lstlisting}

\textbf{Escala} — em torno do centro $(c_x, c_y)$: levar ao centro, multiplicar por $s_x$, $s_y$, voltar:
\begin{lstlisting}
x1, y1 = x - cx, y - cy
x2, y2 = x1 * sx, y1 * sy
return (x2 + cx, y2 + cy)
\end{lstlisting}

\textbf{Rotação} — em torno do centro, usando $\cos$ e $\sin$ do ângulo:
\begin{lstlisting}
x_novo = cx + (x-cx)*cos - (y-cy)*sin
y_novo = cy + (x-cx)*sin + (y-cy)*cos
\end{lstlisting}

\section{Modelo 3D}
O arquivo \texttt{transformacoes\_3d.py} usa um cubo (8 vértices). As transformações usam matrizes $4 \times 4$: translação $(t_x, t_y, t_z)$, escala $(s_x, s_y, s_z)$ e rotação em torno do eixo Y. A projeção isométrica converte $(x,y,z)$ em coordenadas 2D para desenho na tela.

\section{Como executar}
\begin{enumerate}
    \item \texttt{pip install pygame-ce}
    \item \texttt{python transformacoes\_2d.py} ou \texttt{python transformacoes\_3d.py}
\end{enumerate}

\end{document}
